\chapter{Heart Block (AV Block)}
\index{Heart Block (AV Block)}
\label{sec:Heart Block (AV Block)}

\section{Overview}
Atrioventricular (AV) block is a conduction disorder in which the electrical impulse from the atria is delayed or completely blocked as it travels through the AV node or His-Purkinje system to the ventricles. It is classified into three degrees of severity: first-degree (PR prolongation), second-degree (intermittent conduction failure), and third-degree (complete heart block with no atrial impulses reaching the ventricles). This condition can range from an asymptomatic ECG finding to a life-threatening emergency requiring pacemaker implantation.
\section{Classification}

\begin{description}[font=\normalfont\bfseries,leftmargin=3em,labelwidth=2.5em]
  \item[Cause] Idiopathic fibrosis, ischemic, or structural
  \item[Organ System] Heart
  \item[Pathophysiology] Impaired conduction through the AV node or His-Purkinje system delaying or blocking electrical impulse transmission from atria to ventricles
  \item[Duration] Chronic (may be acute in ischemic settings)
  \item[Onset] Gradual (can be sudden in acute myocardial infarction)
  \item[Mode of Transmission] Non-communicable
  \item[Clinical Course] Variable; first- and Mobitz I second-degree are generally benign; complete block requires pacing
  \item[Severity] Mild to Severe (depending on degree)
  \item[Extent of Disease] Localized (cardiac conduction system)
  \item[Age Group] Older adults (most common); any age with congenital AV block
  \item[Epidemiology] Prevalence of first-degree AV block ~5\% in adults over 65; complete heart block ~0.04\% in general population
  \item[ICD-11 Code] BC63.2
\end{description}

\section{Causes}
Degenerative fibrosis of the cardiac conduction system (Lenegre disease or Lev disease) is the most common cause of chronic progressive AV block in older adults. Ischemic heart disease, particularly acute myocardial infarction involving the inferior or anterior wall, can cause transient or permanent conduction block. Other causes include cardiomyopathy, myocarditis, cardiac surgery or catheter ablation complications, medication toxicity (beta-blockers, calcium channel blockers, digoxin), and congenital heart defects.

\section{Risk Factors}

\begin{itemize}[noitemsep]
  \item Advanced age (senescence of the conduction system)
  \item Coronary artery disease and acute myocardial infarction
  \item Cardiomyopathy (dilated, hypertrophic, or infiltrative)
  \item Cardiac surgery (valve surgery, septal myectomy, congenital heart repair)
  \item Medications affecting AV conduction (beta-blockers, verapamil, diltiazem, digoxin, amiodarone)
  \item Systemic diseases (sarcoidosis, amyloidosis, Lyme disease, rheumatic fever)
  \item Congenital heart disease (especially corrected transposition of great arteries)
\end{itemize}

\section{Signs and Symptoms}

\subsection{Early symptoms}
\begin{itemize}[noitemsep]
  \item Early manifestations of heart block (av block) may be subtle and nonspecific
\end{itemize}

\subsection{Common symptoms}
\begin{itemize}[noitemsep]
  \item \subsection{Common symptoms}
  \item \textbf{Common symptoms:}
  \item Asymptomatic (common in first-degree and Mobitz I second-degree)
  \item Fatigue and generalized weakness
  \item Dizziness or presyncope
  \item Syncope (Stokes-Adams attacks in complete heart block)
  \item Shortness of breath on exertion
  \item Slow or irregular heartbeat (bradycardia)
  \item Chest pain (if due to myocardial ischemia)
  \item Reduced exercise tolerance
  \item Altered mental status or confusion in severe bradycardia
  \item Pain or discomfort in the affected area
  \item Difficulty performing daily activities
\end{itemize}

\subsection{Less common symptoms}
\begin{itemize}[noitemsep]
  \item Atypical or less common presentations of heart block (av block) may occur
\end{itemize}

\subsection{Emergency warning signs}
\begin{itemize}[noitemsep]
  \item Severe or worsening symptoms of heart block (av block) require immediate medical evaluation
\end{itemize}
\section{Progression and Stages}
First-degree AV block may remain stable for years or slowly progress to higher degrees over many years. Mobitz I second-degree block (Wenckebach) is usually benign and localizes to the AV node. Mobitz II second-degree block often indicates disease below the AV node in the His-Purkinje system and carries an increased risk of progression to complete heart block. Third-degree (complete) AV block results in complete dissociation of atrial and ventricular rhythms, with ventricular escape rhythms that may be unreliable and unstable.

\section{Complications}

\begin{itemize}[noitemsep]
  \item Syncope and falls due to Stokes-Adams attacks
  \item Heart failure from reduced cardiac output in severe bradycardia
  \item Ventricular arrhythmias triggered by bradycardia (torsades de pointes)
  \item Sudden cardiac death (especially in complete heart block with unstable escape rhythm)
  \item Reduced quality of life
  \item Emotional and psychological effects
\end{itemize}

\section{Diagnosis}

\begin{itemize}[noitemsep]
  \item 12-lead ECG demonstrating prolonged PR interval, dropped QRS complexes, or AV dissociation
  \item Holter or event monitoring for intermittent or paroxysmal block
  \item Exercise stress testing to evaluate chronotropic competence
  \item Electrophysiology study (His bundle recording) to localize the level of block
  \item Echocardiogram to assess structural heart disease
  \item Lab tests including electrolytes, thyroid function, and cardiac biomarkers
  \item Imaging tests when structural heart disease is suspected
\end{itemize}

\section{Prevention}

\begin{itemize}[noitemsep]
  \item Manage cardiovascular risk factors (hypertension, diabetes, hyperlipidemia)
  \item Avoid medications that impair AV conduction when possible in at-risk patients
  \item Treat underlying conditions such as ischemic heart disease and cardiomyopathy
  \item Prompt recognition and treatment of Lyme disease and other infectious causes
  \item Regular ECG monitoring in patients on medications that affect conduction
  \item Go for regular check-ups
  \item Follow your doctor's screening recommendations
\end{itemize}

\section{Treatment Options}

\begin{itemize}[noitemsep]
  \item Observation with serial monitoring for asymptomatic first-degree and Mobitz I block
  \item Discontinuation or dose reduction of culprit medications
  \item Atropine or isoproterenol for temporary management of acute symptomatic bradycardia
  \item Temporary transvenous or transcutaneous pacing for acute settings
  \item Permanent pacemaker implantation (single-chamber or dual-chamber) for symptomatic or high-grade block
  \item Cardiac resynchronization therapy (biventricular pacing) when indicated
  \item Lifestyle changes and self-care
  \item Surgery when needed
\end{itemize}

\section{Living with It}

\begin{itemize}[noitemsep]
  \item Follow your treatment plan as prescribed
  \item Keep up with regular appointments
  \item Adjust your daily habits as needed
  \item Reach out to family, friends, or support groups
  \item Keep learning about your condition
\end{itemize}

\section{Outlook}
The outlook for AV block depends on the degree and location of the block: asymptomatic first-degree and Mobitz I block carry an excellent prognosis, while Mobitz II and complete heart block require permanent pacing for symptom relief and prevention of sudden death. With appropriate pacemaker management, most patients have a normal life expectancy and good functional status.
